% ============================================================
%  FN6905 Exotic Options & Structured Products — Final Assignment
%  Combined report: Parts (a), (b), (c)
%
%  Build instructions:
%    1. python generate_figures.py   (saves figures/ subdirectory)
%    2. pdflatex main.tex
%    3. pdflatex main.tex            (second pass for cross-refs)
% ============================================================
\documentclass[12pt,a4paper]{article}

% ── Packages ────────────────────────────────────────────────
\usepackage[a4paper, margin=2.4cm]{geometry}
\usepackage{amsmath,amssymb,amsthm}
\usepackage{booktabs}
\usepackage{graphicx}
\usepackage{float}
\usepackage{xcolor}
\usepackage[hidelinks,colorlinks=true,linkcolor=blue!60!black,
            citecolor=green!50!black,urlcolor=blue!70!black]{hyperref}
\usepackage{caption}
\usepackage{subcaption}
\usepackage{array}
\usepackage{multirow}
\usepackage{enumitem}
\usepackage{microtype}
\usepackage{parskip}
\usepackage{titlesec}
\usepackage{fancyhdr}

% ── Typography ──────────────────────────────────────────────
\setlength{\parskip}{6pt}
\setlength{\parindent}{0pt}

\titleformat{\section}{\large\bfseries\color{blue!30!black}}{\thesection}{1em}{}
\titleformat{\subsection}{\normalsize\bfseries\color{blue!20!black}}{\thesubsection}{1em}{}
\titleformat{\subsubsection}{\normalsize\bfseries}{\thesubsubsection}{1em}{}

% ── Page style ──────────────────────────────────────────────
\pagestyle{fancy}
\fancyhf{}
\lhead{\small FN6905 Exotic Options \& Structured Products --- Final Assignment}
\rhead{\small\thepage}
\renewcommand{\headrulewidth}{0.4pt}

% ── Math shortcuts ──────────────────────────────────────────
\newcommand{\Phi}{\Phi}
\newcommand{\N}[1]{\Phi(#1)}
\DeclareMathOperator{\E}{\mathbb{E}}
\newcommand{\R}{\mathbb{R}}
\newcommand{\dd}{\,\mathrm{d}}
\newcommand{\half}{\tfrac{1}{2}}

% ── Column type ─────────────────────────────────────────────
\newcolumntype{C}[1]{>{\centering\arraybackslash}p{#1}}
\newcolumntype{L}[1]{>{\raggedright\arraybackslash}p{#1}}

% ============================================================
\begin{document}

% ── Title page ──────────────────────────────────────────────
\begin{titlepage}
  \centering
  \vspace*{3cm}
  {\LARGE\bfseries FN6905: Exotic Options \& Structured Products\\[0.6em]
   Final Assignment\par}
  \vspace{1.5cm}
  {\large\itshape Parts (a), (b) and (c): Monte Carlo, Deep PPDE, and PDGM Pricing
   of Barrier and Lookback Options\par}
  \vspace{2cm}
  {\normalsize Nanyang Technological University\\
   Nanyang Business School\par}
  \vspace{0.8cm}
  \begin{center}
    \includegraphics[width=0.5\textwidth]{NTU_Logo.png} \\[1cm]
    \end{center}
  {\large\bfseries Stefanus Gunawan\par}
  \vspace{0.3cm}
  {\normalsize Matriculation Number: G2500310J\par}
  \vspace{0.5cm}
  {\normalsize \today\par}
  \vfill
\end{titlepage}

\tableofcontents
\clearpage

% ============================================================
\section{Part (a): Monte Carlo Pricing}
\label{sec:parta}
% ============================================================

\subsection{Introduction}

This section implements Monte Carlo (MC) pricers for barrier options
(Table~8.1 of the course notes) and floating lookback options
(Propositions~9.1 and~9.5) under the Black-Scholes model.
All prices are compared to their closed-form analytical values.

\subsection{Simulation Methodology}

\subsubsection{Path Generation}

Stock price paths are simulated under the risk-neutral measure using the
exact log-Euler scheme for Geometric Brownian Motion:
\begin{equation}
  S_{t+\Delta t} = S_t \exp\!\Bigl[(r - \tfrac{\sigma^2}{2})\,\Delta t
    + \sigma\sqrt{\Delta t}\,Z\Bigr],
    \quad Z\sim\mathcal{N}(0,1),
  \label{eq:gbm}
\end{equation}
where $r$ is the risk-free rate, $\sigma$ the volatility, and $\Delta t = T/N$
the step size.  Equation~\eqref{eq:gbm} is the exact solution to the GBM SDE;
only the discrete monitoring of barriers and extrema introduces bias.

\subsubsection{Barrier Option Payoffs}

\begin{table}[H]
\centering
\caption{Barrier option payoff conditions ($S_T$ = terminal price,
         $B$ = barrier level).}
\label{tab:barrier_payoffs}
\begin{tabular}{L{3.8cm} C{3.0cm} C{4.5cm}}
\toprule
\textbf{Type} & \textbf{Activation condition} & \textbf{Payoff at $T$} \\
\midrule
Down-and-out call & $\min_t S_t > B$ & $\max(S_T - K, 0)$ \\
Down-and-in call  & $\min_t S_t \le B$ & $\max(S_T - K, 0)$ \\
Up-and-out call   & $\max_t S_t < B$ & $\max(S_T - K, 0)$ \\
Up-and-in call    & $\max_t S_t \ge B$ & $\max(S_T - K, 0)$ \\
Down-and-out put  & $\min_t S_t > B$ & $\max(K - S_T, 0)$ \\
Down-and-in put   & $\min_t S_t \le B$ & $\max(K - S_T, 0)$ \\
Up-and-out put    & $\max_t S_t < B$ & $\max(K - S_T, 0)$ \\
Up-and-in put     & $\max_t S_t \ge B$ & $\max(K - S_T, 0)$ \\
\bottomrule
\end{tabular}
\end{table}

The Monte Carlo estimate is $\hat{V} = e^{-rT}\,\frac{1}{M}\sum_{i=1}^{M}
\text{Payoff}^{(i)}$, with 95\% confidence interval
$\hat{V} \pm 2\,e^{-rT}\,\sigma_{\mathrm{payoff}}/\!\sqrt{M}$.

\subsubsection{Lookback Option Payoffs}

\begin{table}[H]
\centering
\caption{Floating lookback payoffs.}
\label{tab:lookback_payoffs}
\begin{tabular}{L{4cm} C{4cm} C{4cm}}
\toprule
\textbf{Option} & \textbf{Payoff at $T$} & \textbf{Reference} \\
\midrule
Floating PUT  & $M_T - S_T$ & Proposition~9.1 \\
Floating CALL & $S_T - m_T$ & Proposition~9.5 \\
\bottomrule
\end{tabular}
\end{table}

where $M_T = \max_{0\le t\le T}S_t$ and $m_T = \min_{0\le t\le T}S_t$.

\subsection{Closed-Form Expressions}

\subsubsection{Delta Functions}

Define, for $\tau = T - t$:
\begin{equation}
  \delta^\pm_\tau(z)
    = \frac{\ln z + (r \pm \tfrac{\sigma^2}{2})\,\tau}{\sigma\sqrt{\tau}},
  \qquad \delta^-_\tau(z) = \delta^+_\tau(z) - \sigma\sqrt{\tau}.
  \label{eq:delta}
\end{equation}

\subsubsection{Barrier Options (Table 8.1, Reiner-Rubinstein)}

Let $\mu = (r - \sigma^2/2)/\sigma^2$.  Define the building blocks:
\begin{align}
  A &= \eta\Bigl[S\,\Phi\!\bigl(\eta\,\delta^+_T(S/K)\bigr)
        - K e^{-rT}\Phi\!\bigl(\eta\,\delta^-_T(S/K)\bigr)\Bigr], \label{eq:A} \\
  B &= \eta\Bigl[S\,\Phi\!\bigl(\eta\,\delta^+_T(S/B)\bigr)
        - K e^{-rT}\Phi\!\bigl(\eta\,\delta^-_T(S/B)\bigr)\Bigr], \label{eq:B} \\
  C &= \eta\Bigl[\!\left(\frac{B}{S}\right)^{\!2\mu+2}\!
        S\,\Phi\!\bigl(\varphi\,\delta^+_T(B^2/SK)\bigr)
        - K e^{-rT}\!\left(\frac{B}{S}\right)^{\!2\mu}
        \Phi\!\bigl(\varphi\,\delta^-_T(B^2/SK)\bigr)\Bigr], \label{eq:C} \\
  D &= \eta\Bigl[\!\left(\frac{B}{S}\right)^{\!2\mu+2}\!
        S\,\Phi\!\bigl(\varphi\,\delta^+_T(B/S)\bigr)
        - K e^{-rT}\!\left(\frac{B}{S}\right)^{\!2\mu}
        \Phi\!\bigl(\varphi\,\delta^-_T(B/S)\bigr)\Bigr]. \label{eq:D}
\end{align}
Here $\eta = +1$ for calls, $\eta = -1$ for puts;
$\varphi = +1$ for down-barriers, $\varphi = -1$ for up-barriers.
Selected prices: down-and-out call ($B \le K$): $A - C$;
up-and-out put ($B \ge K$): $B - D$.
Knock-in prices follow from parity:
$V_{\mathrm{in}} = V_{\mathrm{vanilla}} - V_{\mathrm{out}}$.

\subsubsection{Proposition 9.1 --- Floating Lookback PUT}

Payoff $= M_T - S_T$.  At $t = 0$, $M_0 = S_0$:
\begin{equation}
  P = M_0 e^{-rT}\Phi\!\bigl(-\delta^-_T(S_0/M_0)\bigr)
    + S_0\Bigl(1 + \frac{\sigma^2}{2r}\Bigr)\Phi\!\bigl(\delta^+_T(S_0/M_0)\bigr)
    - S_0 e^{-rT}\frac{\sigma^2}{2r}\!\left(\frac{M_0}{S_0}\right)^{-2r/\sigma^2}
      \Phi\!\bigl(-\delta^-_T(M_0/S_0)\bigr) - S_0.
  \label{eq:prop91}
\end{equation}

\subsubsection{Proposition 9.5 --- Floating Lookback CALL}

Payoff $= S_T - m_T$.  At $t = 0$, $m_0 = S_0$:
\begin{equation}
  C = S_0\,\Phi\!\bigl(\delta^+_T(S_0/m_0)\bigr)
    - m_0 e^{-rT}\Phi\!\bigl(\delta^-_T(S_0/m_0)\bigr)
    + e^{-rT} S_0 \frac{\sigma^2}{2r}\!\left(\frac{m_0}{S_0}\right)^{-2r/\sigma^2}
      \Phi\!\bigl(\delta^-_T(m_0/S_0)\bigr)
    - S_0\frac{\sigma^2}{2r}\Phi\!\bigl(-\delta^+_T(S_0/m_0)\bigr).
  \label{eq:prop95}
\end{equation}

\subsection{Parameters}

\begin{table}[H]
\centering
\caption{Monte Carlo simulation parameters.}
\label{tab:params_a}
\begin{tabular}{L{3.5cm} C{2.5cm} L{6cm}}
\toprule
\textbf{Parameter} & \textbf{Value} & \textbf{Notes} \\
\midrule
$S_0$ (initial price)    & 100          & \\
$K$ (strike)             & 100          & At-the-money \\
$B$ (barrier)            & 90 / 110     & Down-options / up-options \\
$r$ (risk-free rate)     & 0.05         & \\
$\sigma$ (volatility)    & 0.20         & \\
$T$ (maturity)           & 1.0 year     & \\
$N$ (time steps)         & 252          & Daily monitoring \\
$M$ (paths)              & 100{,}000    & \\
\bottomrule
\end{tabular}
\end{table}

\subsection{Results}

\subsubsection{Barrier Options}

\begin{table}[H]
\centering
\caption{Barrier option prices: Monte Carlo vs.\ closed-form.
  $S_0=100$, $K=100$, $r=0.05$, $\sigma=0.20$, $T=1$, $N=252$, $M=100{,}000$.}
\label{tab:barrier_results}
\begin{tabular}{L{2.8cm} C{1.1cm} C{1.1cm} C{1.7cm} C{1.4cm} C{2.0cm} C{1.4cm} C{1.4cm}}
\toprule
\textbf{Type} & $B$ & \textbf{Cond.} & \textbf{MC} & $\pm2$SE & \textbf{CF} & \textbf{Diff} & \textbf{Diff\%} \\
\midrule
Call down-out & 90  & $B\le K$ & 8.8635  & 0.0923 & 8.6655  & $+0.198$ & $+2.3\%$ \\
Call down-in  & 90  & $B\le K$ & 1.5359  & 0.0338 & 1.7851  & $-0.249$ & $-14.0\%$ \\
Call up-out   & 110 & $B\ge K$ & 0.1564  & 0.0056 & 0.1186  & $+0.038$ & $+31.9\%$ \\
Call up-in    & 110 & $B\ge K$ & 10.2429 & 0.0932 & 10.3320 & $-0.089$ & $-0.9\%$ \\
Put down-out  & 90  & $B\le K$ & 0.1932  & 0.0062 & 0.1512  & $+0.042$ & $+27.8\%$ \\
Put down-in   & 90  & $B\le K$ & 5.3586  & 0.0552 & 5.4223  & $-0.064$ & $-1.2\%$ \\
Put up-out    & 110 & $B\ge K$ & 4.3636  & 0.0529 & 4.0796  & $+0.284$ & $+7.0\%$ \\
Put up-in     & 110 & $B\ge K$ & 1.1882  & 0.0248 & 1.4939  & $-0.306$ & $-20.5\%$ \\
\bottomrule
\end{tabular}
\end{table}

\begin{figure}[H]
  \centering
  \includegraphics[width=0.95\textwidth]{figures/fig1_barrier_mc.png}
  \caption{All eight barrier option prices: Monte Carlo vs.\ closed-form.
           Large errors for knock-out options (up-out/down-out) arise from
           discrete-monitoring bias.}
  \label{fig:barrier_mc}
\end{figure}

The MC prices agree with closed-form to within a few percent for in-the-money
options (knock-in types) but show larger percentage deviations for
out-of-the-money knock-out options (e.g.\ $+31.9\%$ for call up-out).
This is a well-known \emph{discrete-monitoring bias}: the closed-form assumes
continuous barrier monitoring, while simulation checks the barrier only at
$N = 252$ discrete times.  The Broadie-Glasserman-Kou~\cite{bgk1997}
continuity correction $B \to B\,e^{\pm 0.5826\,\sigma\sqrt{\Delta t}}$
reduces this bias to $O(\Delta t)$; we report uncorrected prices for
transparency.

\subsubsection{Lookback Options}

\begin{table}[H]
\centering
\caption{Floating lookback option prices: Monte Carlo vs.\ closed-form.
  $S_0=100$, $r=0.05$, $\sigma=0.20$, $T=1$, $N=252$, $M=100{,}000$.}
\label{tab:lookback_results}
\begin{tabular}{L{4.5cm} C{2.2cm} C{2.2cm} C{1.6cm} C{2cm} C{2cm}}
\toprule
\textbf{Option} & \textbf{Proposition} & \textbf{MC} & $\pm2$SE & \textbf{CF} & \textbf{Diff\%} \\
\midrule
Floating PUT  ($M_T - S_T$) & 9.1 & 13.4291 & 0.0800 & 14.2906 & $-6.0\%$ \\
Floating CALL ($S_T - m_T$) & 9.5 & 16.5870 & 0.0869 & 17.2168 & $-3.7\%$ \\
\bottomrule
\end{tabular}
\end{table}

\begin{figure}[H]
  \centering
  \includegraphics[width=0.60\textwidth]{figures/fig2_lookback_mc.png}
  \caption{Floating lookback option prices (Monte Carlo vs.\ closed-form).
           Both methods agree to within 4--6\%; the gap reflects discrete
           monitoring of the running extremum.}
  \label{fig:lookback_mc}
\end{figure}

MC prices lie below the closed-form for both options because discretely
monitored extrema underestimate the continuous running maximum / overestimate
the minimum.

% ============================================================
\clearpage
\section{Part (b): Deep PPDE Pricing}
\label{sec:partb}
% ============================================================

\subsection{Introduction}

This section applies the \emph{Deep PPDE} framework of
Sabat\'e-Vidales, {\v S}i{\v s}ka, and Szpruch~\cite{sabate2018} to price
lookback and barrier options.  The method uses an LSTM network to encode path
history and frames the pricing problem as a backward stochastic differential
equation (BSDE).

\subsection{Methodology}

\subsubsection{BSDE Representation}

Under the risk-neutral measure, the option price $Y_t$ and its delta
$Z_t = \sigma S_t\,\partial Y_t/\partial S_t$ satisfy:
\begin{equation}
  Y_{t_i} = Y_{t_{i+1}} + r\,Y_{t_{i+1}}\,\Delta t
             - Z_{t_i}^{\top}\Delta W_i,
  \label{eq:bsde}
\end{equation}
where $\Delta W_i = W_{t_{i+1}} - W_{t_i}\sim\mathcal{N}(0, \Delta t)$.
Two LSTM networks approximate $Y_t$ (price) and $Z_t$ (delta) jointly; the
loss is the residual of Eq.~\eqref{eq:bsde} summed over all time steps plus
a terminal loss:
\begin{equation}
  \mathcal{L} = \E\!\left[\sum_{i=0}^{N-1}
    \bigl|Y_{t_i} - Y_{t_{i+1}} - r\,Y_{t_{i+1}}\Delta t
    + Z_{t_i}^\top\Delta W_i\bigr|^2
    + \lambda\,\bigl|Y_T - \text{Payoff}(\{S_t\})\bigr|^2\right].
  \label{eq:ppde_loss}
\end{equation}

\subsubsection{Path-Dependent Extension via Log-Signature}

For path-dependent payoffs, the Deep PPDE computes the
\emph{log-signature} of the price path between consecutive coarse grid
points, which captures the full path geometry as a compact feature vector.
This allows the LSTM to represent the running extremum (lookback) or the
barrier-crossing indicator (barrier) without explicit memory of every tick.

\subsubsection{Parameters}

\begin{table}[H]
\centering
\caption{Parameters for the Deep PPDE experiments.}
\label{tab:params_b}
\begin{tabular}{L{4.5cm} C{3cm} C{3cm}}
\toprule
\textbf{Parameter} & \textbf{Lookback} & \textbf{Barrier} \\
\midrule
$S_0$                        & 1.0       & 1.0 \\
$K$                          & 1.0       & 1.0 \\
$r$                          & 0.05      & 0.05 \\
$\sigma$                     & 0.30      & 0.20 \\
$T$                          & 1.0       & 1.0 \\
$B$ (barrier level)          & ---       & 0.9 / 1.1 \\
Training iterations          & 500       & 3000--5000 \\
Dimensionality $d$           & 1 and 4   & 1 \\
\bottomrule
\end{tabular}
\end{table}

\subsection{Results}

\subsubsection{Lookback Options ($d = 1$, Single Asset)}

\begin{table}[H]
\centering
\caption{Single-asset lookback prices: Deep PPDE vs.\ closed-form vs.\ Monte Carlo.
  $S_0=1$, $r=0.05$, $\sigma=0.30$, $T=1$, 500 training iterations.}
\label{tab:ppde_lookback_1d}
\begin{tabular}{L{4cm} C{1.8cm} C{1.8cm} C{1.8cm} C{2cm} C{2cm}}
\toprule
\textbf{Option} & \textbf{Prop.} & \textbf{Deep PPDE} & \textbf{MC} & \textbf{CF} & \textbf{PPDE vs CF} \\
\midrule
Floating PUT  & 9.1 & 0.2160 & 0.2130 & 0.2330 & $-7.3\%$ \\
Floating CALL & 9.5 & 0.2289 & 0.2242 & 0.2379 & $-3.8\%$ \\
\bottomrule
\end{tabular}
\end{table}

\subsubsection{Lookback Options ($d = 4$, Basket)}

\begin{table}[H]
\centering
\caption{Multi-asset basket lookback ($d = 4$): Deep PPDE vs.\ Monte Carlo.
  No single-asset closed-form applies.}
\label{tab:ppde_lookback_4d}
\begin{tabular}{L{5cm} C{3cm} C{3cm} C{2.5cm}}
\toprule
\textbf{Option} & \textbf{Deep PPDE} & \textbf{MC} & \textbf{Diff\%} \\
\midrule
Basket Floating PUT  & 0.3879 & 0.3711 & $+4.5\%$ \\
Basket Floating CALL & 0.5372 & 0.5288 & $+1.6\%$ \\
\bottomrule
\end{tabular}
\end{table}

The $d = 4$ basket lookback has no analytical formula, demonstrating
that the Deep PPDE scales naturally to multi-dimensional settings where
traditional closed-form methods fail.

\begin{figure}[H]
  \centering
  \includegraphics[width=0.80\textwidth]{figures/fig3_ppde_lookback_loss.png}
  \caption{Deep PPDE training loss (log scale) for lookback PUT and CALL.
    Loss converges within 500 iterations ($\approx 5$ min on CPU).}
  \label{fig:ppde_lookback_loss}
\end{figure}

\subsubsection{Barrier Options}

\begin{table}[H]
\centering
\caption{Barrier option prices: Deep PPDE vs.\ closed-form vs.\ Monte Carlo.
  $S_0=1$, $K=1$, $r=0.05$, $\sigma=0.20$, $T=1$, 3000--5000 iterations.}
\label{tab:ppde_barrier}
\begin{tabular}{L{3cm} C{1.2cm} C{2cm} C{2cm} C{2cm} C{2.2cm}}
\toprule
\textbf{Option} & $B$ & \textbf{Deep PPDE} & \textbf{MC} & \textbf{CF} & \textbf{PPDE vs CF} \\
\midrule
Down-out CALL & 0.9 & 0.0802 & 0.0894 & 0.0867 & $-7.5\%$ \\
Up-out PUT    & 1.1 & 0.0529 & 0.0448 & 0.0408 & $+29.7\%$ \\
Down-in CALL  & 0.9 & 0.0249 & 0.0143 & 0.0179 & $+39.1\%$ \\
Up-in PUT     & 1.1 & 0.0131 & 0.0110 & 0.0149 & $-12.1\%$ \\
\bottomrule
\end{tabular}
\end{table}

\begin{figure}[H]
  \centering
  \includegraphics[width=0.80\textwidth]{figures/fig4_ppde_barrier_loss.png}
  \caption{Deep PPDE training loss for barrier options.  Barrier training is
    harder: the discontinuous knock-out payoff slows convergence, requiring
    3000--5000 iterations.}
  \label{fig:ppde_barrier_loss}
\end{figure}

\subsubsection{In-Out Parity Verification}

\begin{table}[H]
\centering
\caption{In-out parity check: $V_{\mathrm{out}} + V_{\mathrm{in}} \approx V_{\mathrm{vanilla}}$.}
\label{tab:parity}
\begin{tabular}{L{3.5cm} C{2cm} C{2cm} C{2.2cm} C{2.2cm} C{1.8cm}}
\toprule
\textbf{Option} & \textbf{Knock-out} & \textbf{Knock-in} & \textbf{Sum} & \textbf{Vanilla BS} & \textbf{Error} \\
\midrule
Down call ($B=0.9$) & 0.0894 & 0.0143 & 0.1037 & 0.1046 & $0.9\%$ \\
Up put   ($B=1.1$)  & 0.0448 & 0.0110 & 0.0558 & 0.0557 & $0.2\%$ \\
\bottomrule
\end{tabular}
\end{table}

The sum of knock-out and knock-in MC prices closely matches the vanilla
Black-Scholes price, confirming the correctness of the simulation.

\subsection{Discussion}

Both the Deep PPDE and MC prices lie systematically below the closed-form
lookback values.  The closed-form formula assumes \emph{continuous}
barrier monitoring, while both simulation approaches check the extremum
only at $N = 100$ discrete steps; the Broadie-Glasserman-Kou~\cite{bgk1997}
bias $O(\sigma\sqrt{\Delta t})$ explains the discrepancy.

Barrier options converge significantly more slowly than lookback options.
The knock-out condition creates a discontinuity in the price surface: the
option price drops abruptly to zero when the path crosses $B$.
The LSTM interpolates across this discontinuity with difficulty, leading to
erratic loss curves and larger PPDE-vs-CF errors.
The lookback payoff is a continuous function of the running extremum and is
more amenable to BSDE-based neural network training.

% ============================================================
\clearpage
\section{Part (c): PDGM Pricing}
\label{sec:partc}
% ============================================================

\subsection{Introduction}

This section adapts the \emph{Physics-Driven Galerkin Method} (PDGM) of
Zhang et al.~\cite{pdgm2022} to price barrier and lookback options.
The original code prices geometric Asian options using a TensorFlow~1 LSTM;
we re-implement it in PyTorch and replace the Asian payoff with lookback and
barrier terminal conditions.

\subsection{PDGM Architecture}

An LSTM processes the stock-price path $S_0, S_1, \ldots, S_t$ to produce
a hidden state $h_t$ that encodes the path history.  A feedforward network
(FFN) then maps $(S_t,\, t,\, h_t)$ to an option price estimate $f(S_t, t)$.
The total loss is:
\begin{equation}
  \mathcal{L} = \underbrace{\E\!\left[\sum_{i=0}^{N-1}
    \bigl|\partial_t f + r S_t\,\partial_S f
    + \tfrac{1}{2}\sigma^2 S_t^2\,\partial_{SS}f - rf\bigr|^2
  \right]}_{\mathcal{L}_{\mathrm{PDE}}}
  +\; N\cdot\underbrace{\E\!\left[
    \bigl|f(S_T, T;\, h_T) - \text{Payoff}\bigr|^2
  \right]}_{\mathcal{L}_{\mathrm{TC}}},
  \label{eq:pdgm_loss}
\end{equation}
where $\partial_S f$ and $\partial_{SS}f$ are computed via central finite
differences in $S$ with $h_t$ held fixed (the LSTM hidden state is treated
as a fixed encoding of path history when computing spatial gradients).

\subsection{Modifications for Barrier and Lookback Options}

\begin{table}[H]
\centering
\caption{Terminal payoff functions for each option type in the PDGM.}
\label{tab:pdgm_payoffs}
\begin{tabular}{L{3.8cm} L{5.5cm} L{4cm}}
\toprule
\textbf{Option type} & \textbf{Terminal payoff} & \textbf{LSTM role} \\
\midrule
Lookback PUT (Prop.~9.1)  & $M_T - S_T$ & Encodes running maximum \\
Lookback CALL (Prop.~9.5) & $S_T - m_T$ & Encodes running minimum \\
Barrier knock-out         & $\max(S_T{-}K,0)\cdot\mathbf{1}_{\{S \not\text{ crossed } B\}}$ & Encodes barrier breach \\
Barrier knock-in          & $\max(S_T{-}K,0)\cdot\mathbf{1}_{\{S\text{ crossed }B\}}$ & Encodes barrier breach \\
\bottomrule
\end{tabular}
\end{table}

The PDE loss remains the standard Black-Scholes PDE for all option types.
The LSTM naturally tracks the running extremum (lookback) or barrier-crossing
history (barrier) through its hidden state.

\subsection{Hyperparameters}

\begin{table}[H]
\centering
\caption{PDGM network and training hyperparameters.}
\label{tab:pdgm_hyper}
\begin{tabular}{L{5cm} C{3cm} L{4.5cm}}
\toprule
\textbf{Hyperparameter} & \textbf{Value} & \textbf{Notes} \\
\midrule
LSTM hidden size ($n_a$) & 64          & Original: 128 \\
FFN architecture         & $3\times64$ + tanh & Same depth as original \\
Batch size ($M$)         & 256         & \\
Time steps ($N$)         & 50          & $\Delta t = 0.02$ \\
Training epochs          & 2000 / 3000 & Lookback / barrier \\
Learning rate            & $10^{-3}$ (Adam) & Exp.\ decay $\gamma=0.9998$ \\
Gradient clip            & max norm 5.0 & \\
Spatial bump (FD)        & $0.01\,S_t$ & Central differences \\
\bottomrule
\end{tabular}
\end{table}

\subsection{Black-Scholes Parameters}

\begin{table}[H]
\centering
\caption{Black-Scholes parameters for all PDGM experiments.}
\label{tab:pdgm_params}
\begin{tabular}{L{5cm} C{2.5cm} C{2.5cm}}
\toprule
\textbf{Parameter} & \textbf{Symbol} & \textbf{Value} \\
\midrule
Initial stock price   & $S_0$    & 1.0 \\
Risk-free rate        & $r$      & 0.05 \\
Volatility            & $\sigma$ & 0.20 \\
Maturity              & $T$      & 1.0 \\
Strike (barrier)      & $K$      & 1.0 \\
Barrier (down-out call) & $B$    & 0.9 \\
Barrier (up-out put)    & $B$    & 1.1 \\
\bottomrule
\end{tabular}
\end{table}

\subsection{Results}

\subsubsection{Lookback Options}

\begin{table}[H]
\centering
\caption{Lookback option prices: PDGM vs.\ closed-form vs.\ Monte Carlo.
  $S_0=1$, $r=0.05$, $\sigma=0.20$, $T=1$, $N=50$, 2000 epochs.}
\label{tab:pdgm_lookback}
\begin{tabular}{L{3.5cm} C{1.8cm} C{1.8cm} C{1.8cm} C{2.2cm} C{2.2cm}}
\toprule
\textbf{Option} & \textbf{Prop.} & \textbf{CF} & \textbf{PDGM} & \textbf{MC} & \textbf{PDGM vs CF} \\
\midrule
Floating PUT  & 9.1 & 0.1429 & 0.1266 & 0.1245 & $-11.4\%$ \\
Floating CALL & 9.5 & 0.1722 & 0.1642 & 0.1593 & $-4.6\%$ \\
\bottomrule
\end{tabular}
\end{table}

Both PDGM and MC prices lie below the continuous-monitoring closed-form
values due to discrete-monitoring bias.  Notably, the PDGM price is closer
to the closed-form than raw Monte Carlo in both cases: the PDE loss term
pushes the network toward the continuous-time Black-Scholes equation,
partially correcting the discrete-monitoring effect.

\begin{figure}[H]
  \centering
  \includegraphics[width=0.80\textwidth]{figures/fig5_pdgm_lookback_loss.png}
  \caption{PDGM training loss for lookback options (2000 epochs).
    The loss converges smoothly; the PUT converges slightly faster.}
  \label{fig:pdgm_lookback_loss}
\end{figure}

\subsubsection{Barrier Options}

\begin{table}[H]
\centering
\caption{Barrier option prices: PDGM vs.\ closed-form vs.\ Monte Carlo.
  $S_0=1$, $K=1$, $r=0.05$, $\sigma=0.20$, $T=1$, $N=50$, 3000 epochs.}
\label{tab:pdgm_barrier}
\begin{tabular}{L{3cm} C{1.2cm} C{2cm} C{2cm} C{2cm} C{2.2cm} C{2.2cm}}
\toprule
\textbf{Option} & $B$ & \textbf{CF} & \textbf{PDGM} & \textbf{MC} & \textbf{PDGM vs CF} & \textbf{MC vs CF} \\
\midrule
Down-out CALL & 0.9 & 0.0867 & 0.0914 & 0.0923 & $+5.4\%$  & $+6.5\%$ \\
Up-out PUT    & 1.1 & 0.0408 & 0.0398 & 0.0457 & $-2.5\%$  & $+12.0\%$ \\
\bottomrule
\end{tabular}
\end{table}

\begin{figure}[H]
  \centering
  \includegraphics[width=0.80\textwidth]{figures/fig6_pdgm_barrier_loss.png}
  \caption{PDGM training loss for barrier options (3000 epochs).
    The up-out put converges more erratically due to the discontinuous
    up-barrier payoff.}
  \label{fig:pdgm_barrier_loss}
\end{figure}

\subsection{Discussion}

The PDGM achieves reasonable accuracy for both option classes.  The PDE
regularisation term enforces Black-Scholes dynamics along the path, which
helps the model generalise beyond the discrete training paths.  For barrier
options, the PDGM slightly \emph{overshoots} the CF price for the down-out
call ($+5.4\%$) and undershoots slightly for the up-out put ($-2.5\%$).
Both errors are smaller in magnitude than the raw MC bias, suggesting the
PDE constraint provides partial correction.

% ============================================================
\clearpage
\section{Cross-Method Comparison}
\label{sec:comparison}
% ============================================================

\subsection{Barrier Options: All Methods at a Glance}

Table~\ref{tab:cross_barrier} and Figure~\ref{fig:crossmethod} present a
unified comparison of all four pricing methods for the two representative
barrier options.  All results are for $S_0=1$, $K=1$, $r=0.05$, $\sigma=0.20$,
$T=1$.  Part~(a) Monte Carlo prices are re-scaled from $S_0=100$ by dividing
by 100 (linear scaling of GBM).

\begin{table}[H]
\centering
\caption{Cross-method barrier option prices
  ($S_0{=}1$, $K{=}1$, $r{=}0.05$, $\sigma{=}0.20$, $T{=}1$).}
\label{tab:cross_barrier}
\begin{tabular}{L{4cm} C{2cm} C{2.2cm} C{2.2cm} C{2.2cm} C{2.2cm}}
\toprule
\textbf{Option} & \textbf{CF} & \textbf{MC} \newline \textbf{(Part a)} & \textbf{Deep PPDE} \newline \textbf{(Part b)} & \textbf{PDGM} \newline \textbf{(Part c)} & \textbf{Best vs CF} \\
\midrule
Down-out call ($B{=}0.9$) & 0.0867 & 0.0886 & 0.0802 & 0.0914 & PDGM $+5.4\%$ \\
Up-out put    ($B{=}1.1$) & 0.0408 & 0.0436 & 0.0529 & 0.0398 & PDGM $-2.5\%$ \\
\bottomrule
\end{tabular}
\end{table}

\begin{figure}[H]
  \centering
  \includegraphics[width=0.80\textwidth]{figures/fig7_crossmethod.png}
  \caption{Cross-method comparison of barrier option prices.
    All four methods bracket the closed-form value.
    MC and PDGM lie closest to the CF; Deep PPDE shows larger deviations,
    likely due to fewer training iterations.}
  \label{fig:crossmethod}
\end{figure}

\subsection{Lookback Options: Comparison Notes}

A direct cross-method comparison for lookback options is complicated by a
\emph{parameter mismatch}: the Deep PPDE code in Part~(b) uses $\sigma=0.30$
(the default from the original PPDE repository), while Part~(a) MC and
Part~(c) PDGM both use $\sigma=0.20$.  We therefore present lookback results
per-method only; direct numerical comparison would be misleading.

\begin{table}[H]
\centering
\caption{Lookback results per method (note different $\sigma$ for Part b).}
\label{tab:cross_lookback}
\begin{tabular}{L{3cm} C{2cm} C{2cm} C{2cm} C{2cm}}
\toprule
\textbf{Option} & \textbf{Method} & $\sigma$ & \textbf{Price} & \textbf{CF} \\
\midrule
\multirow{3}{*}{Floating PUT}
  & MC (Part a)    & 0.20 & 13.4291 & 14.2906 ($\times$100) \\
  & Deep PPDE (b)  & 0.30 & 0.2160  & 0.2330 \\
  & PDGM (Part c)  & 0.20 & 0.1266  & 0.1429 \\
\midrule
\multirow{3}{*}{Floating CALL}
  & MC (Part a)    & 0.20 & 16.5870 & 17.2168 ($\times$100) \\
  & Deep PPDE (b)  & 0.30 & 0.2289  & 0.2379 \\
  & PDGM (Part c)  & 0.20 & 0.1642  & 0.1722 \\
\bottomrule
\end{tabular}
\end{table}

\subsection{Key Observations}

\begin{enumerate}[leftmargin=*, label=\textbf{\arabic*.}]

\item \textbf{Discrete-monitoring bias} is the dominant source of error for
  both MC and neural-network methods.  All simulation-based prices lie below
  the continuous-monitoring closed-form for lookback options, and all
  methods overestimate the closed-form for the down-out call.

\item \textbf{PDGM vs.\ Deep PPDE}: PDGM achieves smaller errors for barrier
  options in this experiment.  However, PDGM trained for 2000--3000 epochs
  on a single CPU, while the Deep PPDE results use only 500 iterations for
  lookback and 3000--5000 for barriers.  Given more training time, Deep PPDE
  likely converges further.

\item \textbf{Barrier harder than lookback}: Training loss curves
  (Figures~\ref{fig:ppde_barrier_loss} and~\ref{fig:pdgm_barrier_loss})
  show that barrier options are significantly more difficult for neural
  networks.  The discontinuous payoff at $S = B$ creates a sharp edge in
  the price surface that both LSTMs struggle to approximate.

\item \textbf{Scalability}: Only Deep PPDE (Part b) natively handles
  multiple assets; the $d = 4$ basket lookback result shows that the
  BSDE/log-signature approach scales to higher dimensions where closed-form
  formulas do not exist.

\end{enumerate}

% ============================================================
\clearpage
\section*{References}
\addcontentsline{toc}{section}{References}
% ============================================================

\begin{thebibliography}{9}

\bibitem{sabate2018}
M.~Sabat\'e-Vidales, D.~{\v S}i{\v s}ka, and L.~Szpruch.
\newblock Solving path-dependent PDEs with LSTM networks and the log-signature
  transform.
\newblock \emph{arXiv:1806.01313}, 2018.

\bibitem{pdgm2022}
Z.~Zhang, L.~Kang, and Z.~Wu.
\newblock Physics-driven Galerkin method for solving forward and inverse
  problems in pricing geometric Asian options.
\newblock GitHub repository: \texttt{zhaoyu-zhang/PDGM-Geometric\_Asian}, 2022.

\bibitem{bgk1997}
M.~Broadie, P.~Glasserman, and S.~Kou.
\newblock A continuity correction for discrete barrier options.
\newblock \emph{Mathematical Finance}, 7(4):325--349, 1997.

\bibitem{rr1991}
E.~Reiner and M.~Rubinstein.
\newblock Breaking down the barriers.
\newblock \emph{RISK}, 4(8):28--35, 1991.

\bibitem{conze1991}
A.~Conze and R.~Viswanathan.
\newblock Path dependent options: The case of lookback options.
\newblock \emph{Journal of Finance}, 46(5):1893--1907, 1991.

\end{thebibliography}

\end{document}
